% Ad Familiares 1.8 (ad-familiares-01-08)
% Source: https://raw.githubusercontent.com/PerseusDL/canonical-latinLit/master/data/phi0474/phi056/phi0474.phi056.perseus-lat1.xml
% Fetched by scripts/fetch_latin.py

% Dateline (Perseus): Scr. Romae mense Ian. a. 699 (55) .

\ciceroLetterOpener{M. CICERO S. D. P. LENTVLO PROCOS.}

\ciceroSection{1} de omnibus rebus, quae ad te pertinent, quid actum , quid constitutum sit, quid Pompeius susceperit, optime ex M. Plaetorio cognosces, qui non solum interfuit iis rebus, sed etiam praefuit neque ullum officium erga te hominis amantissimi, prudentissimi, diligentissimi praetermisit. ex eodem de toto statu rerum communium cognosces. quae quales sint, non facile est scribere. sunt quidem certe in amicorum nostrorum potestate, atque ita, ut nullam mutationem umquam hac hominum aetate habitura res esse videatur.

\ciceroSection{2} ego quidem, ut debeo et ut tute mihi praecepisti et ut me pietas utilitasque cogit, me ad eius rationes adiungo, quem tu in meis rationibus tibi esse adiungendum putasti sed te non praeterit, quam sit difficile sensum in re publica, praesertim rectum et confirmatum, deponere. verum tamen ipse me conformo ad eius voluntatem, a quo honeste dissentire non possum, neque id facio, ut forsitan quibusdam videar, simulatione; tantum enim animi inductio et me hercule amor erga Pompeium apud me valet, ut, quae illi utilia sunt, et quae ille vult, ea mihi omnia iam et recta et vera videantur; neque, ut ego arbitror, errarent ne adversarii quidem eius, si, cum pares esse non possent, pugnare desisterent.

\ciceroSection{3} me quidem etiam illa res consolatur, quod ego is sum, cui vel maxime concedant omnes, ut vel ea defendam, quae Pompeius velit, vel taceam vel etiam, id quod mihi maxime libet, ad nostra me studia referam litterarum; quod profecto faciam, si mihi per eiusdem amicitiam licebit. quae enim proposita fuerat nobis, cum et honoribus amplissimis et laboribus maximis perfuncti essemus, dignitas in sententiis dicendis, libertas in re publica capessenda, ea sublata totast, nec mihi magis quam omnibus; nam aut adsentiendum est nulla cum gravitate paucis aut frustra dissentiendum.

\ciceroSection{4} haec ego ad te ob eam causam maxime scribo, ut iam de tua quoque ratione meditere. commutata tota ratio est senatus, iudiciorum, rei totius publicae; otium nobis exoptandum est, quod ii, qui potiuntur rerum, praestaturi videntur, si quidam homines patientius eorum potentiam ferre potuerint; dignitatem quidem illam consularem fortis et constantis senatoris nihil est quod cogitemus; amissa culpa est eorum, qui a senatu et ordinem coniunctissimum et hominem clarissimum abalienarunt.

\ciceroSection{5} sed ut ad ea, quae coniunctiora rebus tuis sunt, revertar, Pompeium tibi valde amicum esse cognovi, et eo tu consule, quantum ego perspicio, omnia, quae voles, obtinebis, quibus in rebus me sibi ille adfixum habebit, neque a me ulla res, quae ad te pertineat, neglegetur; neque enim verebor, ne sim ei molestus, cui iucundum erit etiam propter se ipsum, quom me esse gratum videbit.

\ciceroSection{6} tu velim tibi ita persuadeas, nullam rem esse minimam, quae ad te pertineat, quae mihi non carior sit quam meae res omnes; idque cum sentiam, sedulitate mihimet ipse satis facere possum, re quidem ipsa ideo mihi non satis facio, quod nullam partem tuorum meritorum non modo referenda, sed ne cogitanda quidem gratia consequi possum.

\ciceroSection{7} rem te valde bene gessisse rumor erat; exspectabantur litterae tuae, de quibus eramus iam cum Pompeio locuti. quae si erunt allatae, nostrum studium exstabit in conveniendis magistratibus et senatoribus, ceteraque, quae ad te pertinebunt cum etiam plus contenderimus quam possumus, minus tamen faciemus quam debemus.
